% GENERATED FILE. Edit canonical lessons, then run npm run generate:books.
% canonical-source-hash: fnv1a64:75e972a4c2c1f917
% canonical-lessons: SA-C199-vakyam, SA-C199-bhasa, SA-C199-vyakaranam, SA-C199-ganitam, SA-C199-vijnanam

\chapter{Sentence, Language, Grammar, Mathematics, Science}
\label{ch:sa-sentence-language-grammar-mathematics-science}
\hlchaptermodality{199}

\begin{chapteropening}
\textbf{By the end of this chapter:} \emph{I can say a sentence, a language, grammar, mathematics and science.}

Complete the last lesson of chapter 199: I can say a sentence, a language, grammar, mathematics and science.
\end{chapteropening}

\section[\texorpdfstring{\sk{वाक्यम्} (vākyam)}{वाक्यम् (vākyam)}]{\sk{वाक्यम्} (vākyam) — a sentence}

\label{lesson:SA-C199-vakyam}



\begin{quote}
Before the new one: say the Sanskrit for a nest, then the Sanskrit for an egg.
\end{quote}

\subsection*{You'll want to know: \sk{वाक्यम्}}
\textbf{\sk{वाक्यम्}} — \emph{vākyam} — \textquotedblleft{}a sentence\textquotedblright{}.

\begin{grammarlens}[title={What you've built}]
One more word for this chapter.
\end{grammarlens}

\subsection*{Your turn}
\begin{itemize}
\raggedright
  \item \emph{Say it:} \emph{vākyam}
  \item \emph{Say it:} \emph{vākyam}, once more
  \item \emph{Say it:} say \emph{vākyam}
  \item \emph{Recall:} say the Sanskrit for a nest, then the Sanskrit for an egg, then say \emph{vākyam} again
\end{itemize}

\begin{tcolorbox}[breakable,colback=teal!4,colframe=teal!35!black,title={Before you move on}]
What does \sk{वाक्यम्} mean? (\textquotedblleft{}A sentence\textquotedblright{}.) Say it once more.
\end{tcolorbox}
\section[\texorpdfstring{\sk{भाषा} (bhāṣā)}{भाषा (bhāṣā)}]{\sk{भाषा} (bhāṣā) — a language}

\label{lesson:SA-C199-bhasa}



\begin{quote}
Before the new one: say the Sanskrit for an egg, then the Sanskrit for a sentence.
\end{quote}

\subsection*{You'll want to know: \sk{भाषा}}
\textbf{\sk{भाषा}} — \emph{bhāṣā} — \textquotedblleft{}a language\textquotedblright{}.

\begin{grammarlens}[title={What you've built}]
One more word for this chapter.
\end{grammarlens}

\subsection*{Your turn}
\begin{itemize}
\raggedright
  \item \emph{Say it:} \emph{bhāṣā}
  \item \emph{Say it:} \emph{bhāṣā}, once more
  \item \emph{Say it:} say \emph{bhāṣā}
  \item \emph{Recall:} say the Sanskrit for an egg, then the Sanskrit for a sentence, then say \emph{bhāṣā} again
\end{itemize}

\begin{tcolorbox}[breakable,colback=teal!4,colframe=teal!35!black,title={Before you move on}]
What does \sk{भाषा} mean? (\textquotedblleft{}A language\textquotedblright{}.) Say it once more.
\end{tcolorbox}
\section[\texorpdfstring{\sk{व्याकरणम्} (vyākaraṇam)}{व्याकरणम् (vyākaraṇam)}]{\sk{व्याकरणम्} (vyākaraṇam) — grammar}

\label{lesson:SA-C199-vyakaranam}



\begin{quote}
Before the new one: say the Sanskrit for a sentence, then the Sanskrit for a language.
\end{quote}

\subsection*{You'll want to know: \sk{व्याकरणम्}}
\textbf{\sk{व्याकरणम्}} — \emph{vyākaraṇam} — \textquotedblleft{}grammar\textquotedblright{}.

\begin{grammarlens}[title={What you've built}]
One more word for this chapter.
\end{grammarlens}

\subsection*{Your turn}
\begin{itemize}
\raggedright
  \item \emph{Say it:} \emph{vyākaraṇam}
  \item \emph{Say it:} \emph{vyākaraṇam}, once more
  \item \emph{Say it:} say \emph{vyākaraṇam}
  \item \emph{Recall:} say the Sanskrit for a sentence, then the Sanskrit for a language, then say \emph{vyākaraṇam} again
\end{itemize}

\begin{tcolorbox}[breakable,colback=teal!4,colframe=teal!35!black,title={Before you move on}]
What does \sk{व्याकरणम्} mean? (\textquotedblleft{}Grammar\textquotedblright{}.) Say it once more.
\end{tcolorbox}
\section[\texorpdfstring{\sk{गणितम्} (gaṇitam)}{गणितम् (gaṇitam)}]{\sk{गणितम्} (gaṇitam) — mathematics}

\label{lesson:SA-C199-ganitam}



\begin{quote}
Before the new one: say the Sanskrit for a language, then the Sanskrit for grammar.
\end{quote}

\subsection*{You'll want to know: \sk{गणितम्}}
\textbf{\sk{गणितम्}} — \emph{gaṇitam} — \textquotedblleft{}mathematics\textquotedblright{}.

\begin{grammarlens}[title={What you've built}]
One more word for this chapter.
\end{grammarlens}

\subsection*{Your turn}
\begin{itemize}
\raggedright
  \item \emph{Say it:} \emph{gaṇitam}
  \item \emph{Say it:} \emph{gaṇitam}, once more
  \item \emph{Say it:} say \emph{gaṇitam}
  \item \emph{Recall:} say the Sanskrit for a language, then the Sanskrit for grammar, then say \emph{gaṇitam} again
\end{itemize}

\begin{tcolorbox}[breakable,colback=teal!4,colframe=teal!35!black,title={Before you move on}]
What does \sk{गणितम्} mean? (\textquotedblleft{}Mathematics\textquotedblright{}.) Say it once more.
\end{tcolorbox}
\section[\texorpdfstring{\sk{विज्ञानम्} (vijñānam)}{विज्ञानम् (vijñānam)}]{\sk{विज्ञानम्} (vijñānam) — science}

\label{lesson:SA-C199-vijnanam}



\begin{quote}
Before the new one: say the Sanskrit for grammar, then the Sanskrit for mathematics.
\end{quote}

\subsection*{You'll want to know: \sk{विज्ञानम्}}
\textbf{\sk{विज्ञानम्}} — \emph{vijñānam} — \textquotedblleft{}science\textquotedblright{}.

\begin{grammarlens}[title={What you've built}]
One more word for this chapter.
\end{grammarlens}

\subsection*{Your turn}
\begin{itemize}
\raggedright
  \item \emph{Say it:} \emph{vijñānam}
  \item \emph{Say it:} \emph{vijñānam}, once more
  \item \emph{Say it:} say \emph{vijñānam}
  \item \emph{Recall:} say the Sanskrit for grammar, then the Sanskrit for mathematics, then say \emph{vijñānam} again
\end{itemize}

\begin{tcolorbox}[breakable,colback=teal!4,colframe=teal!35!black,title={Before you move on}]
What does \sk{विज्ञानम्} mean? (\textquotedblleft{}Science\textquotedblright{}.) Say it once more.
\end{tcolorbox}
